\chapter{Casos de Uso Iniciais}
\label{cap:casos-de-uso}

\begin{objetivos}
\begin{itemize}
  \item Reproduzir, no seu contexto, os quatro roteiros pr{\'a}ticos apresentados (planejamento, avalia{\c c}{\~a}o, feedback e sala de aula ativa);
  \item Adaptar os exemplos de \emph{prompts} {\`a} sua {\'a}rea do conhecimento e {\`a} sua disciplina;
  \item Aplicar o checklist de valida{\c c}{\~a}o docente antes de disponibilizar materiais gerados com IA.
\end{itemize}
\end{objetivos}

\section{Introdu{\c c}{\~a}o aos Cen{\'a}rios Pr{\'a}ticos}

Este cap{\'i}tulo foi desenhado para servir como um guia de aplica{\c c}{\~a}o r{\'a}pida. Nele, apresentamos quatro casos de uso t{\'i}picos enfrentados por professores do ensino superior e da educa{\c c}{\~a}o profissional, detalhando as entradas do docente, as chamadas de agentes e os artefatos finais resultantes. A Tabela~\ref{tab:casos-indice} permite localizar rapidamente o caso mais pr{\'o}ximo da sua necessidade.

\begin{table}[htbp]
  \centering
  \small
  \caption{Como Escolher o Caso de Uso Adequado}
  \label{tab:casos-indice}
  \begin{tabularx}{\linewidth}{p{3.2cm}X}
    \toprule
    \textbf{Necessidade Docente} & \textbf{Caso e {\'A}rea Ilustrativa} \\
    \midrule
    Estruturar uma disciplina inteira, definir n{\'i}veis AIAS e cronograma semestral & Caso~1 -- \emph{Estruturas de Dados e Algoritmos} (Ci{\^e}ncia da Computa{\c c}{\~a}o / Engenharias) \\
    \addlinespace
    Elaborar uma avalia{\c c}{\~a}o formativa completa (enunciado, rubrica e quest{\~o}es objetivas) & Caso~2 -- \emph{Banco de Dados} (Ci{\^e}ncia da Computa{\c c}{\~a}o / Engenharias) \\
    \addlinespace
    Gerar feedback formativo inclusivo para estudantes com dislexia ou TDAH & Caso~3 -- Qualquer {\'a}rea com entrega escrita ou c{\'o}digo-fonte \\
    \addlinespace
    Transformar uma aula expositiva em atividade ativa de \emph{Peer Instruction} & Caso~4 -- \emph{Termodin{\^a}mica} (F{\'i}sica / Engenharias) \\
    \bottomrule
  \end{tabularx}
\end{table}

Cada caso segue uma estrutura fixa -- \emph{Contexto e Desafio}, \emph{Execu{\c c}{\~a}o no OpenCode} e \emph{Artefato Entregue} -- facilitando a leitura transversal e a adapta{\c c}{\~a}o para outras disciplinas.

\begin{atencaopedagogica}
\textbf{Como invocar skills e subagentes}: O OpenCode adota duas sintaxes distintas de acionamento. \textbf{Skills} s{\~a}o invocadas com a barra no campo de mensagem, na forma \texttt{/ia-educacao-rubrica} (ou \texttt{/aias-consultant}); alternativamente, podem ser mencionadas no texto do \emph{prompt} pelo slug. \textbf{Subagentes} s{\~a}o acionados com o arroba, na forma \texttt{@planejador-pedagogico} (ou \texttt{@coach-de-feedback-formativo}), ou pela sele{\c c}{\~a}o no menu suspenso da interface. Nos exemplos deste cap{\'i}tulo, os \emph{prompts} seguem essas conven{\c c}{\~o}es: \texttt{@} para o subagente e slug simples (ou \texttt{/} quando isolado) para a skill. Os conceitos de skill e de subagente, bem como os cat{\'a}logos completos, encontram-se respectivamente nos Cap{\'i}tulos~\ref{cap:skills-opencode} e~\ref{cap:agentes-skills}.
\end{atencaopedagogica}

\section{Caso 1: Planejando uma Disciplina do Zero com \emph{Backward Design}}

\subsection{Contexto e Desafio}
O professor assume a disciplina de \emph{Estruturas de Dados e Algoritmos} (60 horas) e deseja estruturar um plano de ensino inovador, alinhado ao \emph{Referencial do MEC} \cite{brasil2026referencial}, com metodologias ativas e enquadramento intencional na Escala AIAS.

\subsection{Execu{\c c}{\~a}o no OpenCode}
No terminal ou na interface gr{\'a}fica, o professor invoca o subagente consultivo:

\begin{lstlisting}[language=bash]
opencode
\end{lstlisting}

\begin{tcolorbox}[colback=white, colframe=boxborder, arc=1mm, title={Comando Enviado pelo Professor}]
\small
``Atue como o subagente \path{@planejador-pedagogico}. Analise a ementa de \emph{Estruturas de Dados} (listas, pilhas, filas, {\'a}rvores bin{\'a}rias, grafos e algoritmos de ordena{\c c}{\~a}o). Proponha um planejamento reverso de 15 semanas com tr{\^e}s momentos avaliativos formativos e especifique o n{\'i}vel AIAS recomendado para cada unidade tem{\'a}tica.''
\end{tcolorbox}

No OpenCode, a forma can{\^o}nica de acionar o subagente {\'e} a men{\c c}{\~a}o pelo nome no pr{\'o}prio texto do \emph{prompt}, na sintaxe ``\path{@planejador-pedagogico}'', ou a sele{\c c}{\~a}o do agente no menu suspenso da interface gr{\'a}fica.

\subsection{Artefato Entregue}
O subagente aplica os frameworks (AIAS, Bloom e \emph{Backward Design}) e indica as skills do reposit{\'o}rio que apoiam a implementa{\c c}{\~a}o do plano (incluindo \path{aias-consultant}, \path{ia-educacao-avaliacao}, \path{ia-educacao-rubrica} e \path{ia-educacao-planejamento-didatico}), gerando:
\begin{itemize}
  \item \textbf{Resultados Desejados}: Compet{\^e}ncias centrais focadas nos n{\'i}veis \emph{Analisar} e \emph{Criar} (avaliar complexidade assint{\'o}tica e projetar estruturas adequadas para problemas reais);
  \item \textbf{Enquadramento AIAS Progressivo}:
    \begin{itemize}
      \item \emph{Semanas 1 a 4 (Ponteiros e Listas)}: N{\'i}vel AIAS~1 (Sem IA), garantindo fixa{\c c}{\~a}o manual de refer{\^e}ncias de mem{\'o}ria;
      \item \emph{Semanas 5 a 10 ({\'A}rvores e Algoritmos)}: N{\'i}vel AIAS~2 (Planejamento Assistido) e AIAS~3 (Colabora{\c c}{\~a}o na depura{\c c}{\~a}o de casos de teste);
      \item \emph{Semanas 11 a 15 (Projeto com Grafos)}: N{\'i}vel AIAS~4 (IA Integral), onde a equipe projeta um sistema de rotas usando IA para co-gerar implementa{\c c}{\~o}es sob rigorosa auditoria discente.
    \end{itemize}
  \item \textbf{Plano de Aulas}: Cronograma de 15 semanas articulado com desafios pr{\'a}ticos em laborat{\'o}rio.
\end{itemize}

\section{Caso 2: Elaborando uma Avalia{\c c}{\~a}o Formativa com Rubrica e MCQs}

\subsection{Contexto e Desafio}
Para o fechamento da segunda unidade de uma disciplina de \emph{Banco de Dados}, o docente precisa de um instrumento avaliativo aut{\^e}ntico com declara{\c c}{\~a}o formal AIAS e uma rubrica anal{\'i}tica detalhada.

\subsection{Execu{\c c}{\~a}o no OpenCode}
\begin{tcolorbox}[colback=white, colframe=boxborder, arc=1mm, title={Comando Enviado pelo Professor}]
\small
``Atue como o subagente \path{@construtor-de-avaliacoes}. Crie uma avalia{\c c}{\~a}o formativa sobre \emph{Modelagem Relacional e Normaliza{\c c}{\~a}o (1FN a 3FN)} para um cen{\'a}rio de com{\'e}rcio eletr{\^o}nico. A atividade ser{\'a} realizada no n{\'i}vel AIAS~2. Em seguida, elabore a rubrica anal{\'i}tica em 4 n{\'i}veis com a skill \path{ia-educacao-rubrica} e adicione duas quest{\~o}es de m{\'u}ltipla escolha diagn{\'o}sticas com distratores baseados em equ{\'i}vocos conceituais t{\'i}picos usando a skill \path{ia-educacao-mcq}.''
\end{tcolorbox}

\subsection{Artefato Entregue}
O subagente gera um arquivo Markdown/LaTeX contendo:
\begin{enumerate}
  \item \textbf{Cabe{\c c}alho Institucional com AIAS}: T{\'i}tulo, instru{\c c}{\~o}es de entrega e a declara{\c c}{\~a}o expl{\'i}cita: ``\emph{Esta atividade adota o N{\'i}vel AIAS 2 -- Planejamento Assistido por IA. Voc{\^e} pode debater ideias de tabelas com a ferramenta, mas os diagramas l{\'o}gicos e as senten{\c c}as DDL devem ser integralmente redigidos por voc{\^e}}'';
  \item \textbf{Desafio Aut{\^e}ntico}: Um cen{\'a}rio real de um e-commerce com anomalias de inser{\c c}{\~a}o e atualiza{\c c}{\~a}o propositais para o discente normalizar;
  \item \textbf{Rubrica Anal{\'i}tica em 4 N{\'i}veis}:
    \begin{itemize}
      \item Crit{\'e}rio 1: Identifica{\c c}{\~a}o de Chaves e Depend{\^e}ncias Funcionais;
      \item Crit{\'e}rio 2: Aplica{\c c}{\~a}o das Formas Normais (1FN a 3FN);
      \item Crit{\'e}rio 3: Integridade Referencial e Restri{\c c}{\~o}es;
      \item Crit{\'e}rio 4: Transpar{\^e}ncia no Uso de IA (coer{\^e}ncia com AIAS~2).
    \end{itemize}
  \item \textbf{Dois Itens MCQs Diagn{\'o}sticos}: Cada distrator com uma nota explicativa indicando exatamente a falha de racioc{\'i}nio do aluno que escolheu aquela alternativa (e.g., confundir depend{\^e}ncia parcial com transitiva).
\end{enumerate}

\section{Caso 3: Ciclo de Feedback Formativo e Inclus{\~a}o}

\subsection{Contexto e Desafio}
O professor recebeu a entrega de um estudante de gradua{\c c}{\~a}o com diagn{\'o}stico de dislexia. O trabalho apresenta excelentes solu{\c c}{\~o}es algor{\'i}tmicas, mas o relat{\'o}rio textual possui estrutura{\c c}{\~a}o truncada e falta a an{\'a}lise de casos de teste de borda.

\subsection{Execu{\c c}{\~a}o no OpenCode}
\begin{tcolorbox}[colback=white, colframe=boxborder, arc=1mm, title={Comando Enviado pelo Professor}]
\small
``Atue como o subagente \path{@coach-de-feedback-formativo}. Analise a submiss{\~a}o anexa em rela{\c c}{\~a}o {\`a} nossa rubrica. O estudante possui laudo de dislexia. Elabore um feedback formativo nas tr{\^e}s dimens{\~o}es (\emph{Feed Up}, \emph{Feed Back}, \emph{Feed Forward}) aplicando as diretrizes de DUA da skill \path{ia-educacao-dua}: use frases curtas em ordem direta, listas pontuadas, limite as a{\c c}{\~o}es de melhoria a duas prioridades e adicione uma pergunta ativadora da metacogni{\c c}{\~a}o.''
\end{tcolorbox}

\subsection{Artefato Entregue}
O subagente gera um parecer acolhedor, objetivo e isento de sobrecarga textual:
\begin{itemize}
  \item \textbf{Feed Up}: Relembra a meta principal (implementar o algoritmo de Dijkstra com fila de prioridades em complexidade $O(E \log V)$);
  \item \textbf{Feed Back}: Destaca em primeiro lugar o acerto do c{\'o}digo (estrutura de dados correta e boa modulariza{\c c}{\~a}o); aponta em seguida a aus{\^e}ncia de testes com grafos desconexos;
  \item \textbf{Feed Forward}: Duas micro-a{\c c}{\~o}es imediatas (passo 1: adicionar uma condi{\c c}{\~a}o de parada para v{\'e}rtices inalcan{\c c}{\~a}veis; passo 2: testar uma entrada com 3 v{\'e}rtices isolados);
  \item \textbf{Pergunta Metacognitiva}: ``Como voc{\^e} explicaria a l{\'o}gica de parada do algoritmo para um colega sem usar jarg{\~o}es t{\'e}cnicos?''
\end{itemize}

\section{Caso 4: Din{\^a}mica Ativa de Sala Invertida e \emph{Peer Instruction}}

\subsection{Contexto e Desafio}
O docente quer transformar uma aula te{\'o}rica expositiva sobre \emph{Termodin{\^a}mica (Segunda Lei)} em uma sess{\~a}o ativa presencial de 100 minutos baseada em discuss{\~a}o por pares.

\subsection{Execu{\c c}{\~a}o no OpenCode}
\begin{tcolorbox}[colback=white, colframe=boxborder, arc=1mm, title={Comando Enviado pelo Professor}]
\small
``Utilize as skills \path{ia-educacao-sala-invertida} e \path{ia-educacao-peer-instruction} para estruturar a aula de Segunda Lei da Termodin{\^a}mica. Preciso do roteiro de leitura pr{\'e}-aula com tr{\^e}s perguntas disparadoras de d{\'u}vida e tr{\^e}s \emph{ConcepTests} conceituais para vota{\c c}{\~a}o presencial com flashcards/celular.''
\end{tcolorbox}

\subsection{Artefato Entregue}
O OpenCode estrutura:
\begin{itemize}
  \item \textbf{Guia Pr{\'e}-Aula (15 minutos discentes)}: Leitura dirigida de duas p{\'a}ginas com perguntas de autorreflex{\~a}o para o aluno trazer anotadas;
  \item \textbf{Din{\^a}mica Presencial}: Divis{\~a}o dos 100 minutos em mini-blocos explicativos de 10 minutos intercalados por \emph{ConcepTests};
  \item \textbf{Quest{\~o}es Conceituais}: Itens qualitativos sem f{\'o}rmulas num{\'e}ricas que for{\c c}am os estudantes a discutir em duplas sobre entropia e irreversibilidade antes da segunda vota{\c c}{\~a}o.
\end{itemize}

\section{Cen{\'a}rios por {\'A}rea STHEM}

As quatro disciplinas ilustrativas dos casos anteriores concentram-se nas {\'a}reas de tecnologia e ci{\^e}ncias exatas. Para evidenciar que o acervo atende a todo o espectro STHEM (Science, Technology, Humanities, Engineering, Mathematics), esta se{\c c}{\~a}o apresenta um cen{\'a}rio adicional por {\'a}rea, com a skill que o OpenCode seleciona espontaneamente ao receber o \emph{prompt}. Os cinco cen{\'a}rios foram validados no OpenCode real: em todos eles, o agente identificou corretamente as skills do acervo a partir da descri{\c c}{\~a}o da tarefa, sem necessidade de citar o nome da skill no \emph{prompt} -- demonstrando que a detec{\c c}{\~a}o sem{\^a}ntica funciona a partir do objetivo pedag{\'o}gico.

\begin{dicaopencode}
\textbf{Valida{\c c}{\~a}o emp{\'i}rica dos roteiros deste cap{\'i}tulo}: os quatro casos de uso (Casos~1 a~4) e os cinco cen{\'a}rios STHEM desta se{\c c}{\~a}o foram executados no OpenCode real (vers{\~a}o 1.18.29, modo headless \texttt{opencode run}). Em todos, as skills citadas foram efetivamente carregadas pela ferramenta \texttt{skill()} e os artefatos esperados foram produzidos (planos, rubricas com r{\'o}tulos can{\^o}nicos, feedbacks tridimensionais e ConcepTests). Um detalhe t{\'e}cnico observado: ao delegar via \texttt{@subagente}, o OpenCode atual exige um identificador de modelo v{\'a}lido no frontmatter do agente (ex.: \texttt{model: opencode/glm-5.3-flash}); o valor \texttt{model: any} do acervo -- um marcador de independ{\^e}ncia de provedor -- n{\~a}o {\'e} interpretado como ID e faz a delega{\c c}{\~a}o direta falhar. Nesse caso, o agente principal assume o papel do subagente e executa as mesmas skills, preservando o resultado; para restaurar a delega{\c c}{\~a}o via Task, basta substituir \texttt{model: any} por um ID de modelo real no arquivo do agente em \texttt{.opencode/agents/}.
\end{dicaopencode}

Os perfis por {\'a}rea seguem a tabela de adapta{\c c}{\~o}es embutida nos subagentes do reposit{\'o}rio, que orientam n{\'i}veis AIAS t{\'i}picos e formatos preferenciais de avalia{\c c}{\~a}o para cada {\'a}rea (Tabela~\ref{tab:sthem-cenarios}).

\begin{table}[htbp]
  \centering
  \small
  \caption{Cen{\'a}rios Validados por {\'A}rea STHEM}
  \label{tab:sthem-cenarios}
  \begin{tabularx}{\linewidth}{p{2.2cm}p{4.6cm}p{3.6cm}X}
    \toprule
    \textbf{{\'A}rea} & \textbf{Cen{\'a}rio Docente} & \textbf{Skills Selecionadas (validado)} & \textbf{N{\'i}vel AIAS T{\'i}pico} \\
    \midrule
    Science & Relat{\'o}rio de laborat{\'o}rio de microbiologia com an{\'a}lise de dados (Biologia) & \path{ia-educacao-escrita}, \path{ia-educacao-avaliacao-projeto} & 1 (prova), 3 (relat{\'o}rio), 4 (dados) \\
    \addlinespace
    Technology & Projeto em equipe de API com IA integral e auditoria discente (Ci. da Computa{\c c}{\~a}o) & \path{ia-educacao-pbl}, \path{ia-educacao-avaliacao-projeto} & 3 (c{\'o}digo assistido), 4 (projeto), 5 (fronteiras) \\
    \addlinespace
    Humanities & Quest{\~o}es discursivas de an{\'a}lise de fontes prim{\'a}rias (Hist{\'o}ria) & \path{ia-educacao-avaliacao} & 2 (pesquisa), 3 (escrita), 5 (co-cria{\c c}{\~a}o) \\
    \addlinespace
    Engineering & Rubrica de memorial de c{\'a}lculo com crit{\'e}rio de transpar{\^e}ncia de IA (Eng. Civil) & \path{ia-educacao-rubrica} & 1 (c{\'a}lculo), 3 (projeto), 4 (otimiza{\c c}{\~a}o) \\
    \addlinespace
    Mathematics & Banco de quest{\~o}es de limites com distratores por equ{\'i}voco (C{\'a}lculo 1) & \path{ia-educacao-banco-questoes}, \path{ia-educacao-mcq} & 1 (dom{\'i}nio), 2 (explora{\c c}{\~a}o), 3 (verifica{\c c}{\~a}o) \\
    \bottomrule
  \end{tabularx}
\end{table}

\subsection{Science: Relat{\'o}rio de Laborat{\'o}rio (Biologia)}

\begin{tcolorbox}[colback=white, colframe=boxborder, arc=1mm, title={Prompt Validado -- Science}]
\small
``Sou professor de Biologia e preciso elaborar o roteiro de um relat{\'o}rio de laborat{\'o}rio de microbiologia que avalie m{\'e}todo cient{\'i}fico e an{\'a}lise de dados, com n{\'i}vel AIAS 3.''
\end{tcolorbox}

\textbf{Resposta do OpenCode (validada)}: o agente seleciona espontaneamente \path{ia-educacao-escrita} (roteiro de relat{\'o}rio t{\'e}cnico, metodologia com IA como apoio sem perder a voz autoral) e \path{ia-educacao-avaliacao-projeto} (rubrica com crit{\'e}rios de an{\'a}lise de dados e marcos de entrega), sinalizando ainda que a declara{\c c}{\~a}o do n{\'i}vel AIAS~3 deve ser calibrada com \path{aias-consultant}.

\subsection{Technology: Projeto Colaborativo com IA Integral (Ci{\^e}ncia da Computa{\c c}{\~a}o)}

\begin{tcolorbox}[colback=white, colframe=boxborder, arc=1mm, title={Prompt Validado -- Technology}]
\small
``Quero planejar uma atividade de projeto onde equipes desenvolvem uma API com assist{\^e}ncia integral de IA (n{\'i}vel AIAS 4), com entregas incrementais e auditoria discente do c{\'o}digo gerado.''
\end{tcolorbox}

\textbf{Resposta do OpenCode (validada)}: o agente seleciona \path{ia-educacao-pbl} (din{\^a}mica de equipes com entregas incrementais) e \path{ia-educacao-avaliacao-projeto} (rubricas e auditoria de projeto com IA, coerentes com o n{\'i}vel AIAS~4).

\subsection{Humanities: An{\'a}lise de Fontes Prim{\'a}rias (Hist{\'o}ria)}

\begin{tcolorbox}[colback=white, colframe=boxborder, arc=1mm, title={Prompt Validado -- Humanities}]
\small
``Elabore tr{\^e}s quest{\~o}es discursivas sobre a Revolu{\c c}{\~a}o Industrial que avaliem an{\'a}lise de fontes prim{\'a}rias.''
\end{tcolorbox}

\textbf{Resposta do OpenCode (validada)}: o agente seleciona \path{ia-educacao-avaliacao} -- e distingue corretamente que, por se tratarem de quest{\~o}es \emph{discursivas} (e n{\~a}o de m{\'u}ltipla escolha), a skill \path{ia-educacao-mcq} n{\~a}o {\'e} a adequada. A arg{\'u}cia de sele{\c c}{\~a}o demonstra que as descri{\c c}{\~o}es das skills guiam a escolha com precis{\~a}o.

\subsection{Engineering: Rubrica de Memorial de C{\'a}lculo (Engenharia Civil)}

\begin{tcolorbox}[colback=white, colframe=boxborder, arc=1mm, title={Prompt Validado -- Engineering}]
\small
``Crie uma rubrica anal{\'i}tica de 4 n{\'i}veis para avaliar o memorial de c{\'a}lculo de estruturas de concreto armado, incluindo crit{\'e}rio de transpar{\^e}ncia no uso de IA (n{\'i}vel AIAS 1 para o c{\'a}lculo).''
\end{tcolorbox}

\textbf{Resposta do OpenCode (validada)}: o agente seleciona \path{ia-educacao-rubrica} -- rubricas anal{\'i}ticas em n{\'i}veis, alinhadas a Bloom, AIAS e DUA, com o crit{\'e}rio de transpar{\^e}ncia no uso de IA. O cen{\'a}rio evidencia a independ{\^e}ncia entre Bloom e AIAS: o c{\'a}lculo pode ser AIAS~1 enquanto a an{\'a}lise cr{\'i}tica do memorial admite n{\'i}veis superiores.

\subsection{Mathematics: Banco de Quest{\~o}es com Distratores (C{\'a}lculo 1)}

\begin{tcolorbox}[colback=white, colframe=boxborder, arc=1mm, title={Prompt Validado -- Mathematics}]
\small
``Monte um banco de 15 quest{\~o}es sobre limites, categorizadas por n{\'i}vel de Bloom, com distratores baseados em equ{\'i}vocos conceituais t{\'i}picos (ex.: confundir limite com valor da fun{\c c}{\~a}o).''
\end{tcolorbox}

\textbf{Resposta do OpenCode (validada)}: o agente seleciona \path{ia-educacao-banco-questoes} (estrutura do banco e blueprint por n{\'i}vel de Bloom) e \path{ia-educacao-mcq} (distratores plaus{\'i}veis ancorados em equ{\'i}vocos reais de racioc{\'i}nio).

\begin{atencaopedagogica}
\textbf{Como reproduzir os cen{\'a}rios}: os cinco \emph{prompts} acima foram testados no OpenCode real e selecionam corretamente as skills sem cit{\'a}-las pelo nome -- basta descrever a tarefa com clareza. O docente pode adaptar o tema (ex.: trocar ``Revolu{\c c}{\~a}o Industrial'' por ``Modernismo brasileiro'') mantendo a estrutura do pedido; a sele{\c c}{\~a}o sem{\^a}ntica da skill permanece a mesma. Para maior determinismo, a skill pode ainda ser citada explicitamente (\texttt{/slug} ou men{\c c}{\~a}o no texto).
\end{atencaopedagogica}

\section{Checklist de Valida{\c c}{\~a}o Docente Pr{\'e}-Aplica{\c c}{\~a}o}

Antes de disponibilizar qualquer material gerado com o aux{\'i}lio do OpenCode aos estudantes, o professor deve realizar a checagem dos cinco pontos da Tabela~\ref{tab:checklist-docente}.

\begin{table}[htbp]
  \centering
  \small
  \caption{Checklist de Valida{\c c}{\~a}o Docente}
  \label{tab:checklist-docente}
  \begin{tabularx}{\linewidth}{clp{3.5cm}X}
    \toprule
    \textbf{Item} & \textbf{Crit{\'e}rio de Controle} & \textbf{Pergunta de Checagem} & \textbf{A{\c c}{\~a}o em Caso de N{\~a}o} \\
    \midrule
    \textbf{1} & Coer{\^e}ncia Pedag{\'o}gica & A tarefa espelha os objetivos reais do curso? & Ajuste os verbos com a skill \path{ia-educacao-bloom}. \\
    \addlinespace
    \textbf{2} & Declara{\c c}{\~a}o AIAS & O n{\'i}vel de IA est{\'a} expl{\'i}cito no enunciado? & Adicione a cl{\'a}usula de integridade com \path{aias-consultant}. \\
    \addlinespace
    \textbf{3} & Auditoria Factual & As f{\'o}rmulas e c{\'o}digos foram rodados? & Execute o teste real e audite com \path{ia-educacao-verificacao}. \\
    \addlinespace
    \textbf{4} & Acessibilidade DUA & A linguagem {\'e} clara e inclusiva? & Simplifique par{\'a}grafos densos com \path{ia-educacao-dua}. \\
    \addlinespace
    \textbf{5} & Autonomia Humana & A avalia{\c c}{\~a}o final depende do professor? & Garanta que o parecer passe pelo seu crivo cr{\'i}tico. \\
    \bottomrule
  \end{tabularx}
\end{table}
